\documentclass[11pt,a4paper]{article}
\usepackage[utf8]{inputenc}
\usepackage[T1]{fontenc}
\usepackage[spanish]{babel}
\usepackage{amsmath,amssymb}
\usepackage{booktabs}
\usepackage{geometry}
\usepackage{hyperref}
\usepackage{enumitem}
\usepackage{array}
\usepackage{listings}
\usepackage{microtype}
\geometry{margin=2.5cm}

\title{\textbf{Propuesta de Evolución Gradual del Algoritmo Sargon}\\
Aprendizaje de pesos, valor contextual de piezas y profundidad adaptativa}
\author{Nota de trabajo}
\date{\today}

\lstset{
  basicstyle=\ttfamily\small,
  frame=single,
  breaklines=true,
  columns=fullflexible,
  showstringspaces=false
}

\begin{document}
\maketitle

\begin{abstract}
Esta nota propone una evolución incremental y controlable del algoritmo Sargon mediante aprendizaje a partir de partidas propias y partidas expertas en formato PGN. La idea rectora es preservar en todo momento una línea base estable: cuando no existe experiencia acumulada, el comportamiento del sistema debe coincidir exactamente con el Sargon clásico. A partir de esa base se introducen, de manera independiente y medible, tres mecanismos de aprendizaje: ajuste de los pesos de la función de valor, estimación contextual del valor de las piezas y profundidad adaptativa de búsqueda. Se propone además un esquema de autoridad de las fuentes de experiencia, un mecanismo de confianza gradual y una estrategia experimental que permita verificar el aporte de cada componente antes de avanzar hacia modelos más complejos.
\end{abstract}

\section{Principio rector}
El objetivo no es reemplazar Sargon por una arquitectura completamente distinta, sino construir una evolución progresiva del algoritmo existente.

La condición fundamental es:
\[
\boxed{\text{sin experiencia} \Rightarrow \text{Sargon clásico}}
\]

y, a medida que se acumula evidencia:
\[
\boxed{
\text{Sargon adaptativo}
=
\text{Sargon clásico}
+
\text{correcciones aprendidas}
}
\]

En ausencia de experiencia:
\[
\Delta w = 0,\qquad
\Delta v = 0,\qquad
\Delta d = 0,
\]
por lo que:
\[
\boxed{
Sargon_{\text{adaptativo}}
=
Sargon_{\text{clásico}}
}
\]

Esta propiedad debe verificarse mediante pruebas automáticas.

\section{Arquitectura conceptual}
Se propone separar el aprendizaje en tres módulos independientes:
\[
\boxed{
\text{Sargon adaptativo}
=
\text{Valor aprendido}
+
\text{Material contextual}
+
\text{Profundidad adaptativa}
}
\]

\section{Módulo I: aprendizaje de pesos de la función de valor}
Supóngase que la evaluación clásica de Sargon puede escribirse como:
\[
V_S(s)=\sum_{i=1}^{n} w_i^{(0)} f_i(s),
\]
donde $f_i(s)$ son características posicionales y $w_i^{(0)}$ los pesos originales.

Entre las características pueden encontrarse material, movilidad, control del centro, estructura de peones, seguridad del rey, desarrollo, espacio, actividad e iniciativa.

La evolución propuesta consiste en mantener el peso original y aprender sólo una corrección:
\[
w_i = w_i^{(0)} + \Delta w_i.
\]

Inicialmente:
\[
\Delta w_i = 0.
\]

La evaluación efectiva será:
\[
V(s)=\sum_i
\left(w_i^{(0)}+\Delta w_i\right)f_i(s).
\]

Esto mantiene el modelo interpretable.

\begin{center}
\begin{tabular}{lcc}
\toprule
Factor & Sargon clásico & Valor aprendido \\
\midrule
Material & 1.00 & 0.98 \\
Movilidad & 0.12 & 0.19 \\
Seguridad del rey & 0.30 & 0.44 \\
Peón aislado & -0.18 & -0.15 \\
\bottomrule
\end{tabular}
\end{center}

\section{Módulo II: valor contextual de las piezas}
Se mantienen como referencia los valores base:
\[
P=1,\qquad N=3,\qquad B=3,\qquad R=5,\qquad Q=9.
\]

Para cada tipo de pieza:
\[
v_p(s)=v_p^{(0)}+\Delta v_p(s).
\]

Por ejemplo, para un caballo:
\[
v_N(s)
=
3+
\alpha_1 C(s)
+
\alpha_2 M(s)
+
\alpha_3 O(s)
+
\alpha_4 F(s),
\]
donde $C$ mide centralización, $M$ movilidad, $O$ calidad de un puesto avanzado y $F$ fase de la partida.

Así, una pieza puede adquirir un valor efectivo dependiente del contexto. Sin experiencia:
\[
\Delta v_p(s)=0.
\]

El mismo principio puede aplicarse a pareja de alfiles, torres en columnas abiertas, dama activa, peones pasados y piezas menores en finales.

\section{Módulo III: profundidad adaptativa}
Sea $d_0$ la profundidad base de Sargon. La profundidad efectiva se define como:
\[
d(s)=d_0+\Delta d(s).
\]

Inicialmente:
\[
\Delta d(s)=0.
\]

Con experiencia, el sistema puede aprender que algunas posiciones merecen más análisis y otras menos.

\begin{center}
\begin{tabular}{lc}
\toprule
Tipo de posición & Profundidad sugerida \\
\midrule
Tranquila & 3 \\
Normal & 4 \\
Compleja & 5 \\
Crítica o táctica & 7 \\
\bottomrule
\end{tabular}
\end{center}

Debe imponerse:
\[
d_{\min}\leq d(s)\leq d_{\max}.
\]

\section{Fuentes de experiencia}
Se consideran dos fuentes principales.

\subsection{Partidas propias}
Las partidas jugadas por el propio motor son válidas, pero su autoridad debe ser moderada:
\[
\omega_{\text{self}}=1.
\]

\subsection{Partidas expertas en PGN}
Las partidas de grandes maestros, motores fuertes o corpus validados reciben mayor autoridad:
\[
\omega_{\text{expert}}>\omega_{\text{self}}.
\]

Por ejemplo:
\[
\omega_{\text{expert}}=5.
\]

Una actualización genérica puede escribirse como:
\[
\Delta \theta
=
\alpha\,
\omega_{\text{source}}\,
\delta\,
\nabla_\theta V.
\]

La procedencia debe conservarse:
\begin{lstlisting}
source=self
source=pgn
source=master
source=engine
\end{lstlisting}

y opcionalmente:
\begin{lstlisting}
authority=1.0
authority=5.0
authority=10.0
\end{lstlisting}

\section{Aprendizaje desde PGN experto}
Sea una posición $s$ donde el experto elige:
\[
a^*.
\]

Sean:
\[
a_1,\ldots,a_n
\]
las jugadas legales.

El PGN aporta evidencia de preferencia:
\[
Q(s,a^*) \gtrsim Q(s,a_j),
\qquad \forall j.
\]

No se asume que $a^*$ sea la única jugada correcta, sino una observación de alta autoridad.

Puede usarse una pérdida de ranking:
\[
L=
\sum_{a\neq a^*}
\max
\left(
0,\,
m-
\left[
V(T(s,a^*))-V(T(s,a))
\right]
\right),
\]
donde $m$ es un margen mínimo.

\section{Aprendizaje desde partidas propias}
El resultado final es:
\[
z=
\begin{cases}
+1 & \text{victoria},\\
0 & \text{tablas},\\
-1 & \text{derrota}.
\end{cases}
\]

Puede aplicarse TD:
\[
\delta_t
=
r_{t+1}
+
\gamma V(s_{t+1})
-
V(s_t).
\]

Con trazas de elegibilidad:
\[
e_t=
\gamma\lambda e_{t-1}
+
\nabla V(s_t),
\]
y:
\[
\theta_{t+1}
=
\theta_t
+
\alpha\delta_t e_t.
\]

Conviene que:
\[
\alpha_{\text{self}}
<
\alpha_{\text{expert}}.
\]

\section{Qué debe aprender cada módulo}
\subsection*{Pesos posicionales}
Los PGN expertos son especialmente útiles porque proporcionan gran cantidad de decisiones:
\[
a^*\succ a.
\]

\subsection*{Valor contextual de piezas}
Debe concentrarse en posiciones informativas: intercambios, sacrificios, simplificaciones, finales, cambios de piezas menores, cambios de dama y entregas de calidad.

\subsection*{Profundidad adaptativa}
El PGN no contiene la profundidad mental del experto, pero puede construirse una señal propia del motor.

Para cada posición experta se ejecuta Sargon a profundidades crecientes:
\[
d=1,2,\ldots,D.
\]

Se define:
\[
d^*
=
\min
\left\{
d:
a_{\text{Sargon}}(s,d)=a_{\text{expert}}
\right\}.
\]

Así se genera un dataset:
\[
(s,d^*).
\]

Este $d^*$ no representa cuánto calculó el maestro, sino cuánto necesita Sargon para reproducir su decisión.

\section{Confianza gradual}
Cada parámetro aprendido debe asociarse con:
\[
(\text{valor aprendido},N,c),
\]
donde $N$ es la evidencia acumulada y $c$ la confianza.

La integración puede formularse como:
\[
w_i^{\text{eff}}
=
(1-c_i)w_i^{(0)}
+
c_iw_i^{(L)},
\]
con:
\[
0\leq c_i\leq1.
\]

Una función sencilla:
\[
c(N)=\frac{N}{N+K}.
\]

Por ejemplo, con $K=1000$:
\[
N=0 \Rightarrow c=0,
\qquad
N=1000 \Rightarrow c=0.5,
\qquad
N=9000 \Rightarrow c=0.9.
\]

La evidencia experta puede incrementar el número efectivo:
\[
N_{\text{efectivo}}
=
\omega_{\text{source}}N.
\]

\section{Evolución por niveles}
\subsection*{Nivel 0: Sargon clásico}
\[
\Delta w=0,\qquad
\Delta v=0,\qquad
\Delta d=0.
\]

\subsection*{Nivel 1: aprendizaje de pesos}
Se habilita únicamente:
\[
w_i^{(0)}\rightarrow w_i.
\]

\subsection*{Nivel 2: valor contextual de piezas}
Se habilita:
\[
v_p(s)
=
v_p^{(0)}
+
\Delta v_p(s).
\]

\subsection*{Nivel 3: profundidad adaptativa}
Se incorpora:
\[
d(s)=d_0+\Delta d(s).
\]

Esta etapa debe introducirse sólo después de demostrar que los niveles anteriores son estables.

\section{Por qué no comenzar con una red neuronal grande}
En las primeras fases conviene utilizar modelos pequeños e interpretables.

Para los pesos:
\[
\Delta w_i
\]
puede ser un conjunto directo de parámetros.

Para material contextual:
\[
\Delta v_p(s)
=
\beta^\top g_p(s).
\]

Para profundidad:
\[
\Delta d(s)
=
\operatorname{round}
\left(
\eta^\top h(s)
\right).
\]

Sólo si una aproximación lineal muestra una limitación objetiva debería reemplazarse por una red neuronal pequeña.

\section{Persistencia}
Se propone una base:
\begin{lstlisting}
sargon_learning.sqlite
\end{lstlisting}

con:
\begin{lstlisting}
learned_weights
piece_context
depth_model
\end{lstlisting}

y estadísticas:
\begin{lstlisting}
source
samples
wins
losses
draws
confidence
last_update
\end{lstlisting}

Además:
\begin{lstlisting}
model_version
created
games_seen
expert_games_seen
\end{lstlisting}

Esto permite auditoría y rollback.

\section{Interfaz de línea de comandos}
Una interfaz posible:
\begin{lstlisting}
sargon --learn-game
sargon --learn-pgn masters.pgn
sargon --learn-pgn masters.pgn --authority 5

sargon --learning off
sargon --learning weights
sargon --learning pieces
sargon --learning full
\end{lstlisting}

\section{Métricas obligatorias}
Deben conservarse simultáneamente:
\[
V_S(s)
\]
y:
\[
V_L(s).
\]

Además:
\begin{lstlisting}
move_baseline
move_learned
\end{lstlisting}

Una métrica importante es:
\[
P
\left(
a_{\text{learned}}
\neq
a_{\text{Sargon}}
\right).
\]

También deben medirse:

\begin{itemize}
    \item coincidencia con jugadas expertas;
    \item error de valoración;
    \item estabilidad de pesos;
    \item variación por fuente;
    \item profundidad media;
    \item nodos buscados;
    \item resultados de partidas;
    \item diferencia entre evaluación clásica y aprendida.
\end{itemize}

\section{Arquitectura resumida}
\begin{verbatim}
                    +----------------------+
                    |   Sargon baseline    |
                    | pesos originales     |
                    | piezas clasicas      |
                    | depth original       |
                    +----------+-----------+
                               |
                               v
                    +----------------------+
                    | learned corrections  |
                    |                      |
                    | Delta pesos          |
                    | Delta piezas         |
                    | Delta profundidad    |
                    +----------+-----------+
                               |
                        confianza c(N)
                               |
                               v
                    +----------------------+
                    |  Sargon adaptativo   |
                    +----------------------+
\end{verbatim}

Fuentes:
\begin{verbatim}
partidas propias --- autoridad baja ----+
                                         +--> aprendizaje
PGN expertos ------ autoridad alta -----+
\end{verbatim}

La integración conceptual:
\[
\boxed{
\theta_{\text{eff}}
=
(1-c)\theta_{\text{Sargon}}
+
c\theta_{\text{learned}}
}
\]

con:
\[
c(0)=0.
\]

\section{Plan experimental recomendado}
La primera implementación debería limitarse al aprendizaje de pesos de evaluación.

La persistencia, sin embargo, debe diseñarse desde el comienzo para admitir:
\begin{itemize}
    \item \texttt{piece\_context};
    \item \texttt{depth\_model}.
\end{itemize}

Antes de avanzar al segundo nivel deberían demostrarse:
\[
\text{PGN experto}
\rightarrow
\text{cambio estable de pesos},
\]

\[
\text{más experiencia}
\rightarrow
\text{mayor confianza},
\]

y:
\[
\text{pesos aprendidos}
\rightarrow
\text{mejor coincidencia con decisiones expertas}.
\]

Sólo entonces debería incorporarse el valor contextual de las piezas.

La profundidad adaptativa debe ser la tercera etapa porque no sólo incorpora conocimiento ajedrecístico, sino también una decisión sobre cuánto recurso computacional asignar.

\section{Conclusión}
La propuesta organiza la evolución de Sargon alrededor de tres capacidades:

\begin{enumerate}
    \item aprender cuánto pesan los factores posicionales;
    \item aprender cuánto vale una pieza según el contexto;
    \item aprender cuánto análisis merece una posición.
\end{enumerate}

Las tres comparten una propiedad esencial:
\[
\boxed{
\text{experiencia cero}
\Rightarrow
\text{Sargon clásico}
}
\]

y sólo modifican el comportamiento cuando existe evidencia suficiente.

Esto permite construir un Sargon progresivamente adaptativo manteniendo interpretabilidad, control experimental, posibilidad de rollback y separación clara entre conocimiento original y conocimiento aprendido.

\end{document}
